\documentclass[10pt,a4paper]{article}
\usepackage{amsmath,amssymb,amsthm,graphicx}
% Class-compatibility preamble for a reader excerpt of
%   Paul Balmer, "Homological support of big objects in tensor-triangulated categories",
%   Journal de l'Ecole polytechnique -- Mathematiques 7 (2020), DOI 10.5802/jep.135.
%
% This file is the source directory's inherited layout adapter, unchanged except for one
% mechanical line: \smaller is declared with \providecommand instead of \newcommand, so that
% the excerpt compiles whether or not the wrapper class already defines \smaller as a
% font-size command. No mathematical macro, symbol, statement, proof step or reference was
% added, changed or removed by this file.
\usepackage{amsmath,amssymb,amsthm,xfrac,xspace,hyperref,graphicx,comment}
\excludecomment{DefTralics}
\providecommand{\C}{}
\providecommand{\bibliofont}{\small}
\newcommand{\ie}{i.e.,\xspace}
\newcommand{\eg}{e.g.,\xspace}
\newcommand{\etc}{etc.\xspace}
\newcommand{\cf}{cf.\xspace}
\newcommand{\resp}{resp.\xspace}
\providecommand{\smaller}{\small}
\newcommand{\Subsection}{\subsection}
\newcommand{\Subsubsection}{\subsubsection}
\newcommand{\skpt}{\smallskip}
\emergencystretch=3em
\providecommand{\loccit}{\textit{loc.\ cit.}}
% Editorial numbering adapter for a reader excerpt of
%   Paul Balmer, "Homological support of big objects in tensor-triangulated categories",
%   Journal de l'Ecole polytechnique -- Mathematiques 7 (2020), DOI 10.5802/jep.135.
%
% The publisher article closes with an appendix whose single section is stepped by hand and
% printed as a capital letter; the theorem it contains therefore prints as A.1. The excerpt
% reproduces that printed form because it retains the appendix theorem alone. This file
% changes no mathematical macro, symbol, statement, proof step or reference.
\renewcommand{\thesection}{\Alph{section}}
\newenvironment{enumeratea}
{\bgroup\def\theenumi{\alph{enumi}}\begin{enumerate}}
{\end{enumerate}\egroup}
\newenvironment{enumeratei}
{\bgroup\def\theenumi{\roman{enumi}}\begin{enumerate}}
{\end{enumerate}\egroup}

\def\hto{\mathrel{\lhook\joinrel\to}}
\def\mto{\mathchoice{\longmapsto}{\mapsto}{\mapsto}{\mapsto}}
\def\To#1{\mathchoice{\xrightarrow{\textstyle\kern4pt#1\kern3pt}}{\stackrel{#1}{\longrightarrow}}{}{}}

\newcommand{\onto}{\mathchoice{\to\hspace*{-5mm}\to}{\to\hspace*{-2.85mm}\to}{\to\hspace*{-3mm}\to}{\to\hspace*{-3mm}\to}}

\newcommand{\into}{\mathchoice{\xymatrix@=.5cm{\ar@{>->}[r]&}}
{\xymatrix@=.3cm{\ar@{>->}[r]&}\hspace*{-1mm}}
{\hspace*{-1mm}\xymatrix@=.25cm{\ar@{>->}[r]&}\hspace*{-1.5mm}}
{\xymatrix@=.1cm{\ar@{>->}[r]&}\hspace*{-1mm}}
}
\usepackage[mathscr]{eucal}
\usepackage{bbm}
\usepackage{cleveref}
\crefdefaultlabelformat{#2\textup{#1}#3}
\usepackage[all,cmtip]{xy}\let\labelstyle\textstyle
\newdir{ >}{{}*!/-10pt/\dir{>}}
\newdir^{ (}{{}*!/-5pt/\dir^{(}}
\newdir_{ (}{{}*!/-5pt/\dir_{(}}

\numberwithin{equation}{section}

\swapnumbers

\newtheorem{Thm}[equation]{Theorem}
\newtheorem*{Thm*}{Theorem}
\newtheorem{Prop}[equation]{Proposition}
\newtheorem{Lem}[equation]{Lemma}
\newtheorem{Cor}[equation]{Corollary}

\theoremstyle{definition}
\newtheorem{Def}[equation]{Definition}
\newtheorem{Exa}[equation]{Example}
\newtheorem{Exas}[equation]{Examples}
\newtheorem{Cons}[equation]{Construction}
\newtheorem{Hyp}[equation]{Hypothesis}
\newtheorem{Rem}[equation]{Remark}
\newtheorem{Rec}[equation]{Recollection}

\DeclareMathOperator{\coker}{coker}
\DeclareMathOperator{\colim}{colim}
\DeclareMathOperator{\cone}{cone}
\DeclareMathOperator{\End}{End}
\DeclareMathOperator{\Hom}{Hom}
\DeclareMathOperator{\id}{id}
\DeclareMathOperator{\Id}{Id}
\DeclareMathOperator{\Ind}{Ind}
\DeclareMathOperator{\Ker}{Ker}
\DeclareMathOperator{\opname}{op}
\DeclareMathOperator{\supp}{supp}
\DeclareMathOperator{\Spc}{Spc}
\DeclareMathOperator{\Spec}{Spec}

\newcommand{\bbZ}{\mathbb{Z}}
\newcommand{\cat}[1]{\mathscr{#1}}
\newcommand{\Endcat}[1]{\End_{\cat #1}}
\newcommand{\Homcat}[1]{\Hom_{\cat #1}}
\newcommand{\inv}{^{-1}}
\newcommand{\isoto}{\overset{\sim}{\longrightarrow}}
\newcommand{\isotoo}{\overset{\sim}{\longrightarrow}}
\newcommand{\op}{^{\opname}}
\let\xto\To
\newcommand{\oto}[1]{\overset{#1} \longrightarrow}
\newcommand{\otoo}[1]{\overset{#1}{\longrightarrow}}
\newcommand{\Paul}[1]{{\color{ForestGreen}#1}}
\newcommand{\SET}[2]{\big\{#1\,\big|\,#2\big\}}
\newcommand{\too}{\mathop{\longrightarrow}\limits}
\newcommand{\unit}{\mathbbm{1}}

\DeclareMathOperator{\Ab}{Ab}
\DeclareMathOperator{\Add}{Add}
\DeclareMathOperator{\coev}{coev}
\DeclareMathOperator{\Der}{D}
\DeclareMathOperator{\Img}{Im}
\DeclareMathOperator{\Loc}{Loc}
\DeclareMathOperator{\modname}{mod}
\DeclareMathOperator{\Mod}{Mod}
\DeclareMathOperator{\SH}{SH}
\newcommand{\smallh}{\mathrm{h}}
\DeclareMathOperator{\Stab}{Stab}
\DeclareMathOperator{\Supp}{Supp}

\newcommand{\ababs}{\textit{ab absurdo}}
\newcommand{\ABp}{\bar{\cA}_{\cB'}}
\newcommand{\AB}{\bar{\cA}_{\cB}}
\newcommand{\AC}{\bar{\cA}_{\cC}}
\newcommand{\adj}{\dashv}
\newcommand{\bat}[1]{\bar{\cat{#1}}}
\newcommand{\bEB}{\bar E_{\cB}}
\newcommand{\bEC}{\bar E_{\cC}}
\newcommand{\boneda}{\bar\yoneda}
\newcommand{\bunit}{\bar\unit}
\newcommand{\cA}{\cat{A}}
\newcommand{\cB}{\cat{B}}
\newcommand{\cC}{\cat{C}}
\newcommand{\cD}{\cat{D}}
\newcommand{\cP}{\cat{P}}
\newcommand{\cS}{\cat{S}}
\newcommand{\cT}{\cat{T}}
\newcommand{\EBp}{E_{\cB'}}
\newcommand{\EB}{E_{\cB}}
\newcommand{\EC}{E_{\cC}}
\newcommand{\fp}{^{\textup{fp}}}
\newcommand{\hBp}{\boneda_{\cB'}}
\newcommand{\hB}{\boneda_{\cB}}
\newcommand{\hC}{\boneda_{\cC}}
\newcommand{\hEB}{\hat E_{\cB}}
\newcommand{\hEC}{\hat E_{\cC}}
\let\hook\hto
\newcommand{\hunit}{\hat\unit}
\newcommand{\ideal}[1]{\langle #1\rangle}
\newcommand{\modd}{\modname\text{\!-}}
\newcommand{\Modd}{\Mod\text{\!-}}
\newcommand{\MS}{\Modd\cat S^c}
\newcommand{\mS}{\modd\cat S^c}
\newcommand{\MT}{\Modd\cat T^c}
\newcommand{\mT}{\modd\cat T^c}
\newcommand{\potimes}[1]{^{\otimes #1}}
\newcommand{\qquadtext}[1]{\qquad\text{#1}\qquad}
\newcommand{\restr}[1]{_{|_{\scriptstyle #1}}}
\newcommand{\sbull}{{\scriptscriptstyle\bullet}}
\newcommand{\SpchS}{\Spch(\cat{S}^c)}
\newcommand{\SpchT}{\Spch(\cat{T}^c)}
\newcommand{\Spch}{\Spc^{\smallh}}
\newcommand{\SpcS}{\Spc(\cat{S}^c)}
\newcommand{\SpcT}{\Spc(\cat{T}^c)}
\newcommand{\tensids}{$\otimes$\nobreakdash-ideals}
\newcommand{\tensid}{$\otimes$\nobreakdash-ideal}
\newcommand{\tensnil}{$\otimes$\nobreakdash-nilpotent}
\newcommand{\vcorrect}[1]{{\vphantom{\vbox to #1em{}}}}
\newcommand{\yoneda}{\mathbbm{h}}
\title{For a big tt-category whose compacts are local, reducedness of the fibre of the comparison map over zero is equivalent to a nilpotence alternative for two composable maps of compacts with vanishing tensor product}
\author{Paul Balmer}
\date{doi:10.5802/jep.135; CC BY 4.0}
\begin{document}
\maketitle
\setcounter{section}{1}
\setcounter{equation}{0}
\noindent\textit{Editorial note.} The counted unit is the author's labelled theorem Thm:hom-local at source lines 881-893 with the proof environment that begins two lines later at source line 895; the pairing is the ordinary adjacent one.\par\medskip
\noindent\textit{Editorial note.} The counted statement is the first numbered item of the excerpt, so with the retained counter values and the appendix numbering adapter it prints as A.1, the printed form of the publisher appendix. This article shares one counter between displayed equations and theorem-class environments, so the context blocks that follow the counted proof are renumbered in the excerpt; their source numbers are recoverable from the labels recorded in the context span descriptions.\par\medskip
\noindent\textit{Editorial note.} The counted proof refers only to its own two items and to the nilpotence theorem of the author's earlier work, which is cited and carried by the retained bibliography; every other label it uses is defined inside the retained statement, proof or context spans, so no reference is left to an external pointer.\par\medskip
\noindent\textit{Editorial note.} The retained preamble is the source directory's inherited layout adapter with the single mechanical change documented in the file itself, followed by an editorial numbering adapter for the appendix and by the author's own packages, theorem declarations and macros; no mathematical macro, symbol, statement or proof step was added, changed or removed.\par\medskip
\noindent\textit{Editorial note.} The retained bibliography is the article's own reference list, flattened at the end of the document so that the citation of the retained proof resolves; no entry was rewritten.\par\medskip
\noindent\textit{Allowed standard prerequisites (editorial).}\begin{itemize}
\item Triangular and homological spectra of a big tt-category. The notions of homological prime, of the comparison map phi to the triangular spectrum, and of the residue tensor functor are recalled in the delivered context. The equivalence proved in the counted theorem is the retained obligation.
\item Local tt-categories and rigidity. Only the standing meaning of local (a tensor product of two non-zero compacts is non-zero) and the standard adjunction tricks using rigidity are assumed; the construction of the Serre tensor-ideal of modules killed by some non-zero compact is retained.
\item Gabriel quotients, Serre tensor-ideals and flatness in locally coherent Grothendieck categories. Standard facts about Gabriel localization and about flat objects in the quotient are assumed; every step of the two implications is retained.
\end{itemize}\medskip
\begin{Thm}
\label{Thm:hom-local}
Let $\cT$ be a big tt-category as in \Cref{Hyp:big-tt} and assume its subcategory of rigid-compact~$\cT^c$ is \emph{local}, meaning that $c\otimes d=0$ with $c,d\in\cT^c$ forces $c=0$ or $d=0$. Then the following are equivalent:
\begin{enumeratei}
\item
\label{it:fiber=*}
The fiber of $\phi\colon \SpchT\onto \SpcT$ over~$(0)$ is reduced to a single point.

\item
\label{it:hom-local}
Whenever two maps~$f$, $g$ in~$\cT^c$ satisfy $f\otimes g=0$ then either $f$ is \tensnil\ on some non-zero compact (that is, there exists $z\in\cT^c$, $z\neq0$ such that $f\potimes{n}\otimes z=0$ for $n\gg1$) or $g$ is \tensnil\ on some non-zero compact.
\end{enumeratei}
\end{Thm}

\begin{proof}
Suppose~\eqref{it:fiber=*} and let us prove~\eqref{it:hom-local}. By standard adjunction tricks using rigidity, we can assume that $f\colon \unit\to x$ and $g\colon \unit\to y$ both start from~$\unit$. Let $\cA=\MT$ as before and consider the subcategory of~$\cA\fp=\mT$ given as follows
\[
\cB_0:=\SET{M\in\cA\fp}{\text{there exists }z\in\cT^c,\,z\neq 0\text{ with }\hat z\otimes M=0}.
\]
Since $\cT^c$ is local, $\cB_0$ is closed under extensions, hence $\cB_0$ is a Serre \tensid\ of~$\cA\fp$. Consider the Gabriel quotient
\[
Q\colon \cA\fp\onto \tilde\cA\fp:=\cA\fp/\cB_0.
\]
Let us write $\tilde z\in\tilde\cA\fp$ for $Q(\hat z)$ for all~$z\in\cT^c$. The maximal \tensids\ of~$\tilde\cA\fp$ are in one-to-one correspondence with those of~$\cA\fp$ containing~$\cB_0$. We claim that these are exactly those in~$\phi\inv(0)$. Indeed, the kernel of the unit $\coev\colon \unit\to \hat z^\vee\otimes \hat z$ in~$\cA\fp$, being killed by~$\hat z$ (in~$\cT^c$ already), becomes zero in~$\tilde\cA\fp$ whenever~$z\neq 0$. In other words, $\tilde\unit$ is a subobject of $\tilde z^\vee\otimes\tilde z$ for every $z\in\cT^c$ non-zero and in particular no proper \tensid\ of~$\tilde\cA\fp$ contains any~$\tilde z$ for $z\neq 0$. Hence the preimage $\phi(\cB)=\yoneda\inv(\cB)=(0)$ is zero for every maximal \tensid\ of~$\cA\fp$ containing~$\cB_0$. Under~\eqref{it:fiber=*}, we just proved that $\tilde\cA\fp$ has a unique maximal Serre \tensid, say~$\cB$. Consider now the kernels $M\into \tilde\unit\xto{\tilde f}\tilde x$ and $N\into \tilde\unit\xto{\tilde g}\tilde y$ in~$\tilde\cA\fp$. Since $\tilde f\otimes M=0$ and $\tilde g\otimes N=0$ by the standard argument (using flatness of~$\tilde x$ and $\tilde y$), we see that $\tilde f$ is \tensnil\ on the Serre \tensid\ $\ideal{M}$ generated by~$M$ and similarly for~$\tilde g$ on~$\ideal{N}$. If, \ababs, both $\ideal{M}$ and $\ideal{N}$ are proper then they are both contained in the unique maximal Serre \tensid~$\cB$ of~$\tilde\cA\fp$. Descending to~$\tilde\cA\fp/\cB$, we therefore have two monomorphisms $\bar f\colon \bunit\into \bar x$ and~$\bar g\colon \bunit\into \bar y$, with $\bar x$ and~$\bar y$ $\otimes$-flat. This implies that $\bar f\otimes \bar g\colon \bunit\into \bar x\otimes\bar y$ is still a monomorphism. This contradicts the assumption that $f\otimes g=0$. Consequently, one of the Serre \tensids~$\ideal{M}$ or~$\ideal{N}$ contains~$\tilde\unit$, hence $\tilde f$ or $\tilde g$ is \tensnil\ in $\tilde\cA\fp=\cA\fp/\cB_0$, and we conclude by definition of~$\cB_0$ that $f$ or $g$ is nilpotent on a non-zero compact.

Conversely, suppose~\eqref{it:hom-local} and let us prove~\eqref{it:fiber=*}. Assume \ababs, that there exist two different homological primes $\cB,\cC\in\phi\inv((0))$ in the fiber of~$\phi$ above zero. Then by \cite[Prop.\,5.3]{Balmer20a}, we have $\EB\otimes \EC=0$ and in particular the tensor of $\eta_\cB\colon \unit\to \EB$ with $\eta_\cC\colon \unit \to \EC$ is zero. It follows from the fact that $\hEB$ and $\hEC$ are colimits of representables and from finite presentation of~$\hunit$ that there exist factorizations $\eta_\cB\colon \unit\xto{f} x\to \EB$ and $\eta_\cC\colon \unit\xto{g} y\to \EC$ with $f\otimes g=0$ and $x,y\in\cT^c$. By~\eqref{it:hom-local}, one of~$f$ or~$g$ is \tensnil\ on some non-zero compact, say~$f$ is. Hence there exists $z\neq0$ in~$\cT^c$ such that $f\potimes{n}\otimes z=0$ for some $n\gg1$. But then $\eta_\cB\potimes{n}\otimes z=0$. Since $\EB$ is a weak ring, $\eta_\cB\otimes\EB$ is a split monomorphism, hence $\eta_\cB\potimes{n}\otimes z=0$ forces $\EB\otimes z=0$. The latter implies $0\neq\hat z\in \cB$, or $\phi(\cB)=\yoneda\inv(\cB)\neq 0$, contradicting the choice of~$\cB$ in the fiber of~$(0)$. So the existence of two points in that fiber is absurd.
\end{proof}

\begin{Hyp}
\label{Hyp:big-tt}
Let $\cT$ be a \emph{`big' tensor-triangulated category}, meaning a rigidly-compactly generated one, as in~\cite{BalmerFavi11}. So~$\cT$ admits small coproducts; its subcategory~$\cT^c$ of compact objects coincides with that of rigid (strongly-dualizable) objects; $\cT^c$ is essentially small and generates~$\cT$ as a localizing subcategory.
\end{Hyp}

\begin{Def}
\label{Def:weak-ring}
We say that an object~$A$ in~$\cT$ with a map~$\eta\colon \unit\to A$ (its `unit') is a \emph{weak ring} if $A\otimes\eta\colon A\to A\otimes A$ is a split monomorphism (whose retraction $A\otimes A\to A$ can be thought of as a unital non-associative multiplication on~$A$).
\end{Def}

\begin{Rec}
\label{Rec:A/C}
Let $\cC$ in~$\cA$ be a localizing subcategory (\ie closed under coproducts, extensions, subobjects and quotients). We have a Gabriel quotient~$\cA/\cC$
\begin{equation}
\label{eq:A/C}
\vcenter{\xymatrix{
\cat{A} \ar@<-.3em>@{->>}[d]_-{Q_\cC}
\\
\cA/\cat{C} {}^{\vcorrect{.6}} \ar@<-.3em>@{_{ (}->}[u]_-{R_\cC}
}}
\end{equation}
where $Q_\cC$ is the universal exact functor with kernel~$\cC$, its right adjoint $R_\cC$ is fully-faithful and $Q_\cC\circ R_\cC\cong\Id_{\cA/\cC}$. We shall only consider subcategories $\cC$ that are $\otimes$-ideal ($\cA\otimes \cC\subseteq\cC$), in which case $\cA/\cC$ inherits a unique tensor structure such that $Q_\cC$ is a tensor functor. When $\cC$ is clear from the context, we often write $\bar{X}$ instead of $Q_\cC(\hat X)$ for $X\in\cT$. We also write $\bat{A}_{\cC}$ (or just $\bat{A}$) for $\cA/\cC$ and
\begin{equation}
\label{eq:hC}
\boneda_{\cC}\colon \cT\too \bat{A}_\cC=\cat{A}/\cC
\end{equation}
for the composite $Q_\cC\circ \yoneda$, a coproduct-preserving cohomological tensor-functor.
\end{Rec}

\begin{Prop}
\label{Prop:A/C}
With notation as in~\Cref{Rec:A/C}, we have:
\begin{enumeratea}
\item
\label{it:A/C-a}
The Gabriel quotient $\bat{A}=\cA/\cC$ is a Grothendieck category, that is closed monoidal. Its tensor product preserves colimits in each variable.
\item
\label{it:A/C-b}
Every $X\in\cT$ has $\otimes$-flat image $\bar X=Q_\cC(\hat X)$ in~$\bat{A}$. If $x\in \cT^c$ then $\bar x$ is rigid.
\item
\label{it:A/C-c}
Every object of~$\bat{A}$ is a quotient of a coproduct $\bigsqcup_{i\in I}\,\bar x_i$ with all $x_i\in \cT^c$.

\item
\label{it:A/C-d}
Every injective object of~$\bat{A}$ is of the form~$\bar E$ for some pure-injective~$E$ in~$\cT$, uniquely characterized by the property~$\hat E\cong R_\cC(\bar E)$ in~$\cA$. In particular, $\bar E$ is $\otimes$-flat. Furthermore, $\boneda_\cC\colon \cT\to \bat{A}$ induces, for every $X\in\cT$, an isomorphism
\[
\Homcat{T}(X,E)\isoto \Homcat{{\bat{A}}}(\bar X,\bar E).
\]
\item
\label{it:A/C-e}
The right adjoint $R_\cC\colon \bat{A}\to \cat{A}$ is lax-monoidal and closed; more precisely for every $M\in\cA$ and $N\in\bat{A}$, if we denote by $[-,-]$ the internal-homs, we have
\[
[M,R_\cC(N)]_{\cA}\cong R_{\cC}([Q_\cC(M),N]_{\bat{A}}).
\]
\item
\label{it:A/C-f}
For every injective $\bar E$ in~$\bat{A}$, the internal-hom $[-,\bar E]_{\bat{A}}\colon \bat{A}\op\to \bat{A}$ is exact.
\end{enumeratea}
\end{Prop}

\begin{proof}
These claims follow from \cite[App.\,A]{BalmerKrauseStevenson20} and \cite[\S2\,\&\,App.\,A]{BalmerKrauseStevenson19}. Let us prove~\eqref{it:A/C-f}, for~\cite[Lem.\,2.8]{BalmerKrauseStevenson19} is only stated for~$\cA$. We have
\[
[-,\bar E]_{\bat{A}}\circ Q_{\cC}=[Q_{\cC}(-),\bar E]_{\bat{A}}\cong Q_{\cC}R_{\cC}([Q_\cC(-),\bar E]_{\bat{A}})\underset{\eqref{it:A/C-e}}{\cong} Q_{\cC}\circ[-,R_\cC(\bar E)]_{\cat{A}}
\]
and the latter is exact by~\cite[Lem.\,2.8]{BalmerKrauseStevenson19}, since $R_\cC\bar E$ is injective in~$\cA$ itself. As~pre\-composition with $Q_\cC$ detects exactness for functors on~$\bat{A}$, we get the result.
\end{proof}

\begin{Cons}
\label{Cons:E_C}
For $\cC\subseteq\cA$ localizing \tensid, the injective hull of~$\bar{\unit}$ in~$\AC$ comes from~$\hC$, by \Cref{Prop:A/C}\,\eqref{it:A/C-d}. So there exists a morphism
\[
\eta_\cC\colon \unit\to \EC
\]
in~$\cT$ such that $\bar\eta_\cC=\hC(\eta_\cC)\colon \bar{\unit}\into \bar E_\cC$ is the injective hull. It is characterized by $R_\cC(\bEC)\cong\hEC$ in~$\cA$. Since $\bEC$ is flat and injective, $\bEC\otimes\bar\eta_{\cC}$ is a split monomorphism in~$\AC$, hence by \Cref{Prop:A/C}\,\eqref{it:A/C-d} again, the retraction exists in~$\cT$, meaning that $(\EC,\eta_\cC)$ is a \emph{weak ring} in~$\cT$ in the sense of \Cref{Def:weak-ring}.
\end{Cons}

\begin{Rec}
\label{Rec:A/B}
A particular class of localizing $\otimes$-ideals $\cC\subseteq\cA$ are those generated by Serre $\otimes$-ideals $\cB\subseteq\cA\fp=\mT$ of finitely presented objects (where \tensid\ only means $\cA\fp\otimes\cB\subseteq\cB$ of course). Explicitly $\cC=\Loc(\cB)$ is the smallest localizing subcategory containing~$\cB$, which is then automatically \tensid\ in~$\cA$. In~some notation, typically in indices, we drop the `$\Loc$' part and write only~$\cB$, like for instance with the canonical cohomological tensor-functor~\eqref{eq:hC}:
\begin{equation}
\label{eq:hB}
\boneda_{\cB}\colon \cT\too \bat{A}_\cB:=\cat{A}/\Loc(\cB).
\end{equation}
Similarly, $\EB$ means $E_{\Loc(\cB)}$ as in \Cref{Cons:E_C} for~$\cC=\Loc(\cB)$.
\end{Rec}

\begin{Prop}
\label{Prop:A/B}
With notation as in~\Cref{Rec:A/B}, we have:
\begin{enumeratea}
\item
\label{it:A/B-a}
The Grothendieck category $\AB$ is still locally coherent. Its finitely presented objects $\bat{A}{}\fp=\cA\fp/\cB$ coincide with the Gabriel quotient of~$\cA\fp$ by~$\cB$.

\item
\label{it:A/B-b}
We can recover $\cB$ and $\Loc(\cB)$ from the pure-injective weak ring~$E_\cB$ of~\Cref{Cons:E_C}, as $\cB=\SET{M\in\cA\fp}{\hEB\otimes M=0}$ and $\Loc(\cB)=\Ker(\hEB\otimes-)$.

\item
\label{it:A/B-c}
Every subobject of a finitely presented object in~$\AB$ is the colimit of its finitely presented subobjects. In particular, if $I\into \bar\unit$ in~$\AB$ then $I=\mathop{\colim}\limits_{\textrm{f.p.\,}M\into I}M$.
\end{enumeratea}
\end{Prop}

\begin{Rec}
\label{Rec:Spch}
The \emph{homological spectrum}~$\SpchT$ consists of all maximal Serre \tensids\ $\cB$ of the abelian subcategory~$\cA\fp=\mT$ of finitely presented $\cT^c$\nobreakdash-modules. We call those~$\cB$ the \emph{homological primes} of~$\cT$. Note that they only depend on the subcategory~$\cT^c$. For each $\cB\in\SpchT$, its preimage in~$\cT^c$ under Yoneda
\[
\phi(\cB)=\yoneda\inv(\cB)\cap \cT^c=\SET{x\in\cT^c}{\hat x\in\cB}
\]
is a (triangular) prime ideal in~$\cT^c$. This defines a surjection~$\phi\colon\SpchT\onto \SpcT$ by~\cite[Cor.\,3.9]{Balmer20a}. Each $\cB\in\SpchT$ yields a coproduct-preserving homological tensor-functor $\hB\colon \cT\to \AB=\cA/\Loc(\cB)$ as in~\eqref{eq:hB}.
\end{Rec}

\begin{Rem}
\label{Rem:simple}
The maximality of~$\cB\in\SpchT$ among the Serre \tensids\ of~$\cA\fp$ tells us that $\AB\fp$ is `simple': it has only the two trivial Serre \tensids, zero and~$\AB\fp$. (See \Cref{Rec:A/B} if necessary.) Simplicity has the following consequence.
\end{Rem}

\begin{Prop}
\label{Prop:ring-in-AB}
Let $\cB\in\SpchT$ and $A\in\AB$ be a weak ring (\Cref{Def:weak-ring}) that is $\otimes$-flat, like $A=\bar E$ for some weak ring~$E\in\cT$. Then either $A=0$ or its unit is a monomorphism $\bunit\into A$.
\end{Prop}

\begin{proof}
This is a basic trick in the proof of the Nilpotence Theorem~\cite{Balmer20a}. Consider the exact sequence $\Ker(\eta)\into \bunit\xto{\eta} A$ in~$\AB$. Since $A$ is flat and $A\otimes\eta$ is a (split) monomorphism, we have $A\otimes\Ker(\eta)=0$. Suppose that $\Ker(\eta)\neq0$ and let us show that $A=0$. By \Cref{Prop:A/B}\,\eqref{it:A/B-c}, we know that $\Ker(\eta)$ is the colimit of its finitely presented subobjects. Take a finitely presented $M\neq0$ with $M\into \Ker(\eta)$. By~flatness of~$A$ and $A\otimes\Ker(\eta)=0$, we see that $A\otimes M=0$. Hence $\Ker(A\otimes-)\cap \AB\fp$ is a non-zero Serre (by flatness) \tensid\ of~$\AB\fp$. As $\AB\fp$ is `simple' (\Cref{Rem:simple}) this forces $\Ker(A\otimes-)\cap \AB\fp=\AB\fp$ to contain~$\bunit$, giving~$A=0$.
\end{proof}

We want to discuss the injectivity of the map $\phi\colon \SpchT\onto \SpcT$, from the homological spectrum (\Cref{Rec:Spch}) to the usual triangular spectrum~$\SpcT$. By a general finite-smashing localization argument, we can reduce to the case where $(0)\in \SpcT$, meaning that $\cT^c$ is \emph{local}. For instance, $\cT$ could be $\SH_{(p)}$ the $p$-local stable homotopy category or $\Der(R)$ the derived category of a commutative local ring~$R$. Under this assumption, we have the following characterization:

\begin{thebibliography}{99}
\bibitem{Balmer05a} P. Balmer - “The spectrum of prime ideals in tensor triangulated categories” , J. reine angew. Math. 588 (2005), p. 149-168
\bibitem{Balmer18a} P. Balmer - “On the surjectivity of the map of spectra associated to a tensor-triangulated functor” , Bull. London Math. Soc. 50 (2018) no. 3, p. 487-495
\bibitem{Balmer19} P. Balmer - “A guide to tensor-triangular classification” , in Handbook of homotopy theory (H. Miller, ed.), Chapman and Hall/CRC, 2019, Available on the author’s web page
\bibitem{Balmer20a} P. Balmer - “Nilpotence theorems via homological residue fields” , Tunis. J. Math. 2 (2020) no. 2, p. 359-378
\bibitem{BalmerCameron20pp} P. Balmer \& J. Cameron - “Computing homological residue fields in algebra and topology” , 2020
\bibitem{BalmerDellAmbrogioSanders16} P. Balmer, I. Dell’Ambrogio \& B. Sanders - “Grothendieck-Neeman duality and the Wirthmüller isomorphism” , Compositio Math. 152 (2016) no. 8, p. 1740-1776
\bibitem{BalmerFavi11} P. Balmer \& G. Favi - “Generalized tensor idempotents and the telescope conjecture” , Proc. London Math. Soc. (3) 102 (2011) no. 6, p. 1161-1185
\bibitem{BensonIyengarKrause08} D. J. Benson, S. B. Iyengar \& H. Krause - “Local cohomology and support for triangulated categories” , Ann. Sci. École Norm. Sup. (4) 41 (2008) no. 4, p. 573-619
\bibitem{BensonIyengarKrause11a} D. J. Benson, S. B. Iyengar \& H. Krause - “Stratifying modular representations of finite groups” , Ann. of Math. (2) 174 (2011) no. 3, p. 1643-1684
\bibitem{BensonIyengarKrause11c} D. J. Benson, S. B. Iyengar \& H. Krause - “Stratifying triangulated categories” , J. Topology 4 (2011) no. 3, p. 641-666
\bibitem{BensonIyengarKrause12b} D. J. Benson, S. B. Iyengar \& H. Krause - “Colocalizing subcategories and cosupport” , J. reine angew. Math. 673 (2012), p. 161-207
\bibitem{BensonIyengarKrause12a} D. J. Benson, S. B. Iyengar \& H. Krause - Representations of finite groups: local cohomology and support , Oberwolfach Seminars , vol. 43 , Birkhäuser/Springer, Basel, 2012
\bibitem{BensonIyengarKrause13} D. J. Benson, S. B. Iyengar \& H. Krause - “Module categories for group algebras over commutative rings” , J. K-Theory 11 (2013) no. 2, p. 297-329, With an appendix by Greg Stevenson
\bibitem{BalmerKrauseStevenson19} P. Balmer, H. Krause \& G. Stevenson - “Tensor-triangular fields: ruminations” , Selecta Math. (N.S.) 25 (2019) no. 1, article ID 13, 36 pages
\bibitem{BalmerKrauseStevenson20} P. Balmer, H. Krause \& G. Stevenson - “The frame of smashing tensor-ideals” , Math. Proc. Cambridge Philos. Soc. 168 (2020) no. 2, p. 323-343
\bibitem{DwyerPalmieri08} W. G. Dwyer \& J. H. Palmieri - “The Bousfield lattice for truncated polynomial algebras” , Homology Homotopy Appl. 10 (2008) no. 1, p. 413-436
\bibitem{HoveyPalmieriStrickland97} M. Hovey, J. H. Palmieri \& N. P. Strickland - Axiomatic stable homotopy theory , Mem. Amer. Math. Soc. , vol. 128, no. 610 , American Mathematical Society, Providence, RI, 1997
\bibitem{HoveyStrickland99} M. Hovey \& N. P. Strickland - Morava $K$-theories and localisation , Mem. Amer. Math. Soc. , vol. 139, no. 666 , American Mathematical Society, Providence, RI, 1999
\bibitem{Krause00} H. Krause - “Smashing subcategories and the telescope conjecture—an algebraic approach” , Invent. Math. 139 (2000) no. 1, p. 99-133
\bibitem{LurieHA} J. Lurie - “Higher algebra” (2017), Online at http://www.math.ias.edu/\textasciitilde{}lurie/
\bibitem{Neeman96} A. Neeman - “The Grothendieck duality theorem via Bousfield’s techniques and Brown representability” , J. Amer. Math. Soc. 9 (1996) no. 1, p. 205-236
\bibitem{Neeman00} A. Neeman - “Oddball Bousfield classes” , Topology 39 (2000) no. 5, p. 931-935
\bibitem{Neeman01} A. Neeman - Triangulated categories , Annals of Math. Studies , vol. 148 , Princeton University Press, Princeton, NJ, 2001
\bibitem{Stevenson13} G. Stevenson - “Support theory via actions of tensor triangulated categories” , J. reine angew. Math. 681 (2013), p. 219-254
\end{thebibliography}
\end{document}
